\documentclass[11pt]{article}
\usepackage[margin=1in]{geometry}
\usepackage[T1]{fontenc}
\usepackage{amsmath,amssymb,amsthm}
\usepackage[hidelinks]{hyperref}
\setlength{\emergencystretch}{2em}
\newtheorem{theorem}{Theorem}
\newtheorem{lemma}[theorem]{Lemma}
\newtheorem{corollary}[theorem]{Corollary}
\newcommand{\A}{\mathbb A}
\newcommand{\C}{\mathbb C}
\newcommand{\F}{\mathbb F}
\newcommand{\Aff}{\operatorname{Aff}}
\newcommand{\Eta}{\mathcal E}
\newcommand{\SL}{\operatorname{SL}}
\title{Inward feedback: low Euler cyclicity, radical towers and refined Mathieu capacities}
\author{Independent internal coupling argument}
\date{16 November 2026}
\begin{document}
\maketitle
\noindent\textbf{Scope and assurance.}
This note combines retained research from different threads. It contains
written arguments, not a registry promotion, independent specialist review
or formal verification. Its conclusions inherit the stated inputs below.
No full Keller, HC4, exceptional Mathieu or higher-dimensional Factorial
conjecture is closed. No agents were messaged. In particular the group
exclusion retains its \emph{specified inertia class}; it is not an
exclusion of every representation with the same abstract group.

\section{A surface restriction obtained by combining topology and exact volume}

We use the retained exact-volume theorem HC4MYML0: a smooth affine complex
surface with a regular nowhere-zero exact two-form cannot admit a
complete-type affine-line ruling in which \emph{every} fibre has a
multiplicity-one component. Equivalently, it contains no open
$\A^1$-bundle over $\mathbb P^1$. This theorem is in
\href{../hc4-triangular-graphs/EXACT_VOLUME_AND_COMPLETE_RULING_COMPONENTS.md}{the canonical exact-volume note}.

Two classical inputs are retained explicitly. Negative logarithmic Kodaira
dimension of a smooth affine complex surface supplies an $\A^1$-fibration
over a smooth curve; see
\href{https://arxiv.org/abs/1408.1328}{Dubouloz--Kishimoto}, Section 1.1.2
and Theorem 2, recalling Miyanishi--Sugie.
Every reduced degenerate fibre is a disjoint union of affine lines,
with positive multiplicities; a ruling with all fibres irreducible and
reduced is a Zariski locally trivial bundle. These are the precise inputs
in
\href{https://arxiv.org/html/2002.12838v1}{Dubouloz}, Section 1,
Proposition 3 and the paragraph after Definition 4.
The source says some results are stated without full detailed proofs;
the same source boundary is retained here.

\begin{lemma}[Orbifold presentation]
For such a ruling $\rho:X\to C$, let the fibre at $c_i$ have components
of multiplicities $m_{ij}$ and put $m_i=\gcd_jm_{ij}$. Then $\pi_1(X)$
is obtained from $\pi_1(C\setminus\{c_i\})$ by imposing
$\gamma_i^{m_i}=1$.
\end{lemma}
\begin{proof}
Away from the finitely many degenerate fibres, the ruling is an
$\A^1$-bundle and its fibre is contractible, so the fundamental groups
of total space and punctured base agree. Inclusion of the complement
of the reduced special fibres into smooth $X$ is surjective on
fundamental groups. Its kernel is normally generated by meridians of
their components. A component of multiplicity $m_{ij}$ projects a
meridian to $\gamma_i^{m_{ij}}$, as follows from the local equation
$\rho-c_i=u^{m_{ij}}$ times a unit at a smooth point of that component.
These relations are equivalent to $\gamma_i^{m_i}=1$.
\end{proof}

\begin{theorem}[Low-Euler cyclicity]\label{thm:cyclic}
Let $X$ be a smooth connected affine complex surface with
$\bar\kappa(X)=-\infty$, finite $\pi_1(X)$ and $\chi(X)\le2$.
If $X$ carries a regular nowhere-zero exact two-form, then
$\pi_1(X)$ is cyclic.
\end{theorem}
\begin{proof}
Choose the ruling just described. The map on fundamental groups
surjects onto $\pi_1(C)$, so finite $\pi_1(X)$ implies
$C=\A^1$ or $\mathbb P^1$.

For $C=\A^1$, the orbifold presentation is a free product of the cyclic
multiple-fibre groups. Such a free product is finite only if at most
one factor is nontrivial, giving the conclusion.

For $C=\mathbb P^1$, Euler integration gives
\[
       \chi(X)=2+\sum_i(s_i-1),
\]
where $s_i$ is the number of reduced components of the special fibre.
Thus $\chi(X)\le2$ forces all fibres irreducible, and $\chi(X)=2$.
Its fundamental group has the spherical-orbifold presentation
\[
 \langle\gamma_1,\ldots,\gamma_r:
       \gamma_1\cdots\gamma_r=1,\ \gamma_i^{m_i}=1\rangle.     \tag{1}
\]
If this finite group is noncyclic, the orbifold is good and spherical.
In particular it has a finite holomorphic cover
$\tau:\mathbb P^1\to\mathbb P^1$ with the exact local orders $m_i$.
This is the standard orientable two-orbifold classification: the bad
sphere orbifolds are a teardrop and unequal-order spindles, whose
fundamental groups are cyclic. See Thurston, Chapter 13, especially
Section 13.2. The finite noncyclic signatures are the spherical triangle
signatures.

Normalize $X\times_{\mathbb P^1,\tau}\mathbb P^1$ to obtain $X'$.
It is connected because the generic $\A^1$ fibre is geometrically
integral. It is finite etale over $X$: locally along an irreducible
fibre $m_iC_i$, the base change $t=s^{m_i}$ and normalization split
$u^{m_i}=s^{m_i}$ into etale branches. Outside these fibres the base
change is etale as well. The normalized ruling $X'\to\mathbb P^1$ has
\emph{reduced} fibres. A fibre may split into several affine lines;
irreducibility of those pullback fibres is not asserted.

The pullback of the given form remains exact and nowhere zero, and
$X'$ is smooth affine because it is finite etale over $X$.
HC4MYML0 now forbids its reduced complete-type ruling. This contradicts
the assumption of noncyclic $\pi_1(X)$.
\end{proof}

The use of \emph{good} orbifold covers is essential. A teardrop has no
manifold orbifold cover; one cannot kill its multiple fibre by an
unavailable global cover. Likewise, a reduced pullback fibre can split:
assuming that it stays irreducible would be incorrect. The stronger
component-selection form of HC4MYML0 is exactly what is needed.

\section{A new entire radical-tower inertia family is excluded}

The knot thread constructs the groups
\[
 H_q=\Aff(W,R)\rtimes S_4,\qquad
 W=\F_2^2,\quad R=\F_2^q,
 \quad |H_q|=24\,2^{3q},                                  \tag{2}
\]
on $W\oplus R$, of degree $D=4\,2^q$.
Here $\Aff(W,R)$ is the additive space of affine functions $W\to R$.
For an edge $e=\{a,b\}$ of the four-point set $W$, let
$s_e$ exchange $a,b$, and let $\delta_e$ be their indicator.
The relevant meridian class consists of
\[
       T_{e,r}(x,z)=(s_ex,z+\delta_e(x)r),\quad r\in R.       \tag{3}
\]
These are involutions, with exactly half the letters moved. They preserve
a common alternating form of rank two. This is the inertia class used
in
\href{../meridional-rank/LIFT_DIMENSION_AND_RADICAL_TOWERS.tex}{the knot radical-tower construction}.

\begin{lemma}\label{lem:deck}
The centralizer of the natural transitive $H_q$ action is the group of
constant radical translations, isomorphic to $R$.
\end{lemma}
\begin{proof}
Constant radical translations are central in $H_q$. Any centralizing
permutation must therefore have the form
$(x,z)\mapsto(f(x),z+c(x))$.
Commuting with every affine-function translation forces $f(x)=x$,
because these functions separate the four base points, for $q>0$.
Commuting with the base $S_4$ then makes $c(x)$ constant.
For $q=0$, the centralizer of the natural $S_4$ action is trivial.
\end{proof}

\begin{theorem}[Decorated radical-tower exclusion]\label{thm:tower}
No actual nonautomorphic complex plane Keller map has natural monodromy
$H_q$ with \emph{every generic branch meridian in the class (3)}
when $q\ge2$. This excludes that entire specified family in degrees
$16,32,64,\ldots$, not merely a finite list of group orders.
\end{theorem}
\begin{proof}
Retain the original finite normalization $N$, original plane source,
and every omitted prime. PF129 gives the exact original ledger
\[
             \sum_C\operatorname{supp}(\sigma_C)+\Eta\le D-1.
                                                               \tag{4}
\]
Every inertia permutation in (3) has support $D/2$. Consequently there
is exactly one branch-image component: two would violate (4), and an
empty branch for a connected finite cover of $\A^2$ would give degree one.

The radical-coset blocks give a genuine degree-four intermediate
$N\to M\to\A^2$, with relative degree $t=2^q$ and base action $S_4$.
Its sole generic branch permutation is a transposition, hence its
quotient moved support is $A=2$.
For its entire etale locus $V$, PF129 yields
\[
                    1\le\chi(V)=4-2-\Eta_M\le2.             \tag{5}
\]
The retained whole-source construction makes $V$ smooth affine,
$\bar\kappa(V)=-\infty$, with the exact nowhere-zero form
$dP\wedge dQ=d(P\,dQ)$.

There is no residual ramification over $V$: (3) fixes every letter
in each fixed base block, and no invisible target component exists.
Thus PF128's residual-support criterion gives that
$Z=h^{-1}(V)\to V$ is finite etale. The whole original plane inside
$Z$ gives $\pi_1(Z)=1$, as in PF128.

The constant radical translations act regularly on each block and
commute with monodromy. They therefore give the full relative deck
group $R$ of the degree-$2^q$ extension. Equivalently use
Lemma~\ref{lem:deck} and the standard generic-deck/monodromy-centralizer
identification. Hence $Z$ is the universal cover of $V$ and
\[
                         \pi_1(V)=R=(C_2)^q.
\]
Theorem~\ref{thm:cyclic} and (5) force this group to be cyclic,
contradicting $q\ge2$.
\end{proof}

No degree-minimality assumption was used. The conclusion does NOT
exclude $H_q$ with other inertia classes: order-four lifts or vertical
radical inertia can carry residual ramification. It does not exclude
all binary affine monodromy or the remaining degree-eight $H_1$ case.
The new knot representations remain valid: knot groups do not carry
the affine-surface exact-volume and Euler hypotheses.

More generally, Theorem~\ref{thm:cyclic} applies to EVERY genuine
intermediate etale locus with $\chi(V)\le2$. If its relative cover is
etale, its relative deck group is cyclic. This adds a low-Euler restriction
to PF128--129 without substituting a plane for an intermediate surface.

\section{Invariant-degree gaps sharpen the Mathieu capacity restriction}

Use the full logarithmic capacity criterion in
\href{../moment-vanishing-audit/LOGARITHMIC_JETS_AND_SPARSE_ALL_RING_SAFETY.md}{the sparse all-ring theorem}.
If $f$ has $k$ distinct left-torus weights and both signs, an eventually
pure $f$ satisfies
\[
    \min(C_+,C_-)<k,\quad
    C_\pm=\max_{\pm i>0}\operatorname{cap}(f_i).             \tag{6}
\]
For a charged component of minimum matrix degree $n$, its leading
right-invariant binary form has maximum root multiplicity $M$, and
\[
 \operatorname{cap}=
 \begin{cases}
 n/2,&M\le1,\\
 n/2-M,&M\ge2.
 \end{cases}                                               \tag{7}
\]
The minimum projective right-orbit lengths are six for $2O$ and twelve
for $2I$, so $M\le\lfloor n/6\rfloor$ or
$\lfloor n/12\rfloor$, respectively.

The classical invariant-degree gaps, also checkable by exact character
averaging, are:
\[
 \C[u,v]^{2O}:\ 8,12,18;\qquad
 \C[u,v]^{2I}:\ 12,20,30.                                  \tag{8}
\]
These are generator degrees; not every even degree occurs.

For $2I$, every charged component not of leading degree twelve has
$n\ge20$. If $20\le n<24$, the root bound makes $M\le1$, so capacity
is at least ten. If $n\ge24$, (7) gives at least
$n/2-\lfloor n/12\rfloor\ge10$.
For $2O$, a charged component of leading degree outside eight and
twelve has $n\ge16$, and its capacity is at least six.

\begin{corollary}[Near-minimal support localization]\label{cor:capacity}
An eventually pure mixed-sign phase with at most ten weights on
$\SL_2/2I$ must have ONE ENTIRE charged side in the degree-twelve
Peter--Weyl component.
For $\SL_2/2O$, an eventually pure mixed-sign phase with at most six
weights must have one entire charged side in the filtered sum of
degrees eight and twelve.
\end{corollary}
\begin{proof}
Otherwise each charged side has a component with capacity at least ten
or six, respectively, contradicting (6).
The invariant-degree gaps mean there are no lower charged pieces
outside the indicated spaces.
\end{proof}

Thus the full mixed-sign family with at most ten weights on $2I$ and
no charged degree-twelve projection is NOT eventually pure, in every
matrix degree. The analogous $2O$ family has at most six weights and
no charged degree-eight/twelve projections. The all-ring purity/safety
classification from the sparse theorem transfers to these enlarged
restricted classes by its same Nullstellensatz and weight-balance
argument: the neutral coefficient ideal plus the product of the two
charged coefficient ideals must be nilpotent. This does not prove the
unrestricted ten- or six-weight statement.

These conditions hold after every allowed automorphism, using the
transformed phase's own degree and support. They combine with strict
leading instability, compact-image interior and finite cyclic stabilizer
constraints. None of those assertions removes the remaining low-degree
charged side by itself.

\section{Already available feedback from the binary Factorial audit}

The completed internal FC(2)/EIT audit credits Long and retains its
unpublished-source assurance boundary. Combined with our exact Gamma
bridge it yields: in four independent real Gaussian coordinates,
a phase invariant under the two independent $O(2)$ block actions is
$f((x_1^2+x_2^2)/2,(x_3^2+x_4^2)/2)$.
Over a characteristic-zero field such a phase is all-marker safe iff
$f=0$; even eventual pure vanishing forces $f=0$.
Indeed pass to a positive power to make purity all-positive, apply
FC(2), and use the polynomial domain property.

Over arbitrary coefficient algebras the corresponding criterion is
nilpotence of the phase coefficient ideal, with normalization its ordinary
coefficient torsion, as already written in the Factorial audit.
This closes the entire two-block radial symmetry sector, in every degree.
It does not imply the same result for arbitrary Gamma shapes or more
blocks. Its homogeneous Haar restriction was already part of the thin
sector closure, so it is not a new full exceptional-Mathieu advance.

\section{Sources and adoption boundary}

The new surface argument uses the two named primary affine-ruling sources,
Thurston's orientable orbifold classification
(\href{https://people.math.osu.edu/davis.12/courses/8800/Thurston13.pdf}{Chapter 13}),
and the retained HC4MYML0 exact-volume theorem.
Its finite-cover normalization and meridian arguments are written above.
PF128--129 and the knot tower are kept at their exact canonical scopes.

For the invariant gaps, the $2I$ presentation is given, for example, in
\href{https://www-users.cse.umn.edu/~webb/Publications/BRSW6.pdf}{Broer--Reiner--Smith--Webb},
Section 4.2.1; the binary-octahedral degrees are classical Klein invariant
data and can independently be checked by the character recurrence
$\chi_n(s)=s\chi_{n-1}(s)-\chi_{n-2}(s)$.
The small character computations certify the required finite gaps only;
the all-degree capacity estimates above are analytic inequalities.

The most consequential new deduction here is
Theorem~\ref{thm:tower}. It should be audited and adopted with its inertia
hypothesis intact. The complete HC4/JC2 and both exceptional MZ cases
remain OPEN. A family exclusion is not a full inverse theorem.
\end{document}
